%% main_v16.tex -- AI & SOCIETY submission (ANONYMISED manuscript, double-blind). Title page = title_page.tex.
%% Bibliography refs_v5.bib. article+natbib here for a compilable PDF; swap to Springer sn-jnl class at final submission.
\documentclass[11pt]{article}
\usepackage[utf8]{inputenc}
\usepackage[T1]{fontenc}
\usepackage{microtype}
\usepackage[margin=1in]{geometry}
\usepackage{setspace}\onehalfspacing
\usepackage{natbib}
\bibliographystyle{plainnat}
\setcitestyle{authoryear,round}
\usepackage{hyperref}\hypersetup{hidelinks}
\title{The limits of audit for self-evaluating AI: authority, correlated error, and the price of judgement}
\author{}\date{}
\begin{document}
\maketitle
\begin{abstract}
When an AI system both produces work and judges it, the worry is that it will game its own evaluation, and the usual fix is structural: separate the producing from the judging, and place a check the optimizer cannot soften outside the loop. That fix is right as far as it goes, and it leaves two problems untouched, both beyond the reach of the audit analogy that inspires it. First, what a deterministic check cannot settle is not unmeasured gaming. Whether a surviving deviation is fraud, error, discovery, or dissent is a normative judgment, so governing it is first a question of who may settle the dispute, and only afterward of measurement. A second problem is mechanical: beyond the incentive divergence that liability already addresses, a machine monitor can fail by sharing a training lineage with the system it checks, wrong in the same places for the same reasons. Human institutions know this failure and answer it with rotation, but shared weights couple two models far more tightly, and whether the coupling can be loosened while frontier systems share their data has barely been asked. Institutions can govern the authority to settle contested deviations, while only engineering can lower the correlation of errors, and the two constrain each other. Because neither can be made clean, our proposal puts a standing price on capture and correlation, repriced by adversarial probes and kept in open contest, and promises no certificate of safety.
\end{abstract}
\noindent\textbf{Keywords:} self-evaluating AI, AI governance, audit, accountability, correlated error, algorithmic insurance

\section{Introduction}

A self-improving system proposes a candidate, scores it, and lets the score pick what to try next. When one process does both the proposing and the scoring, and something pushes it to raise the score, it learns to close the distance between what its evaluation measures and what it was for, closing it from the wrong side. Every organization has met this in the worker who writes the report on his own work, and the response came in two parts. One part splits the unit that produces from the unit that certifies. The other stands up an institution, independent audit, to keep that split honest over time. Recursive AI copies both, as much recent work does, putting a non-gameable monitor outside the optimizing loop and importing audit's guarantees of rule-boundedness, independence, and signed attestation \citep{greenblatt2023,raji2020,manheim2025}.

Can an evaluator be trusted to check a system as capable as itself? Scalable-oversight research already asks that, through debate between models, weak models supervising strong ones, self-critique against a written constitution, and one model judging another \citep{irving2018,burns2023,bai2022,zheng2023}. That work treats the matter as a capability to be trained. We treat it as a governance question, to be handed to someone with standing to decide. Both readings hold, but only the first is being pursued.

Audit history adds a vocabulary of failure the capability literature lacks. Two centuries of corporate self-reporting have shown how a check on a self-interested producer rots. It is captured once the verifier gains a stake in the outcome. It turns to ritual once attestation merely comforts, and to theatre once being audited stands in for being sound. We use that catalogue. But two problems decide whether self-evaluating AI can be governed at all, and the analogy reaches neither.

Take first the cases a fixed rule cannot settle. A deterministic check disposes of the failures it can evaluate and hands the rest to judgement, and the usual reading files that remainder as gaming that slipped the gate. Yet the line between gaming and discovery will not hold still. Specification gaming is defined only against an intended objective the rule failed to capture \citep{krakovna2020,skalse2022}. Read another way, that move is an optimizer finding a path nobody anticipated, or showing the rule was wrong to begin with. Stronger optimizers exploit a misspecified objective sooner, and in sudden jumps as they scale \citep{pan2022}, so the disputed remainder grows with capability. To call a surviving deviation fraud, error, innovation, or dissent is to make a normative judgement, and the live question is who may make it, and for whom.

A second problem is mechanical. A verifier can fail because its interests pull away from the producer's, the failure audit understands best. At Arthur Andersen the consulting fee dwarfed the audit fee and bought the auditor a stake in the client's contentment. A monitor can fail that way too, through the lab that picks it and tunes how readily it vetoes, so liability and licensure stay as pertinent to the monitor's handlers as to any auditor. But a model fails a second way the incentive story misses. Sharing a training lineage with what it checks, it goes blind where the producer goes blind, a correlated error no realignment of interests can repair. What follows works these two problems and how each bounds the answer to the other. Section 2 fixes terms. Section 3 runs a deletion test on the firm analogy. Section 4 develops the contested cases and the triage they demand. Section 5 takes up correlated error and argues that decorrelation is the harder and less portable demand. Section 6 asks who may settle the contested cases, and what gaming surface that authority creates. Section 7 turns diagnosis into a priced and contestable proposal. Section 8 weighs what governance reaches against what only engineering can, prices the capability the cure spends, and marks the limit of what any study here could show.
\section{Co-located generation and evaluation}

This paper is about self-improvement loops in which one optimizing process both produces candidate work, such as code, a plan, or an agent's actions, and scores that work, with the score conditioning the next step. Such a loop is narrower than AI in general and wider than any single product. It sits apart from the external audit of a finished model, where producer and auditor already belong to distinct organizations. It sits apart, too, from benchmark optimization in the abstract, where no single process holds both roles. What marks it out is that \emph{generation} and \emph{evaluation} sit \emph{co-located} in one optimizer, so the parameters or prompts that produced an output also decide whether it passes, under steady pressure to raise the score.

Call a system \emph{unitary} when its governing evaluation can be softened by the optimizing process, as when an agent writes code and also judges that code, and \emph{separated} when at least one governing constraint is fixed beyond the optimizer's reach. The distinction is structural, and the next section starts from it, though it does not finish the matter. A constraint can sit outside the optimizer and still test the wrong property. A verifier can be structurally separate and still fail just as what it checks does. Two gaps follow, one between the form of separation and what the separated check can actually decide, the other between two models that look independent without being so, and Sections 4 and 5 take them up.

\section{Structural separation and its limits}

The structure offers a recombination of familiar parts. Role separation, event logs, provenance, incident closure, and human-machine cooperation each carry a large prior literature, and the claim lies in how they combine. A structural layer separates generation from evaluation, adds a deterministic enforcement floor of model-free checks the optimizer cannot soften, and forces closure with a resource clock outside the loop. An institutional layer borrows audit's three guarantees: a published method the audited party does not set, a verifier with no stake in the outcome, and a signed attestation of what was checked and by whom.

Run a deletion test and the analogy turns out to carry little original mechanism. Strip the firm out, and the prescription mostly survives, since a security engineer reaches a non-gameable monitor from reference-monitor reasoning, a contract theorist reaches the conditional value of monitoring from agency cost, and an auditor supplies the three guarantees as routine. What economic content survives is incomplete contracts, and the multidivisional firm drops away, with the link from incomplete contracts to AI alignment already drawn \citep{hadfieldmenell2019,williamson1985}. The firm-and-audit history, then, works at the level of culture more than mechanism. It hands over the vocabulary of decay named above, and stating its limits in the open fixes how much the analogy may claim before the two harder problems begin.

\section{The contested cases}

Hardening a deterministic floor relocates pressure. The model-free checks catch the failures a fixed rule can evaluate, and what is left moves to judgement-laden reasoning no mechanical predicate settles. Reward-hacking research calls that relocation a whack-a-mole dynamic, where each patched exploit is followed by gaming on a dimension nobody was watching. A natural next step asks how much gaming survives and where it goes, then proposes to measure it. That step assumes the leftover is one well-defined quantity, and once the floor has done its work it is not.

Gaming is not undefinable. Against a stated proxy and a stated true objective, reward hacking has a precise definition, and a proxy can be shown unhackable or not \citep{skalse2022}. The trouble is that the cases a floor cannot reach are the ones where the true objective is unsettled. There, what one reading calls fabricated compliance another calls a workaround, a discovery that the specification was wrong, or a refusal of a norm that deserved refusing. Take the famous reinforcement-learning episodes, where an agent exploits a physics bug or satisfies a reward in an unintended way: these count as failures only against a specification that turned out incomplete \citep{krakovna2020}, and in another frame that behaviour is the system telling its designer the rule was wrong. So the leftover is a set of deviations whose status is contested, and a study that counted the gaming would record which deviations its coders chose to call fraud, lending a count the look of objectivity the matter lacks \citep{porter1995}. The contest is not everywhere of one kind. Some cases settle in time, once the harm or the benefit a deviation produces resolves after the fact what no rule could resolve before it. Some settle by stipulation, once a task fixes the goal and the deviation plainly meets it or fails it. Two kinds resist both: the normatively contested, where reasonable parties disagree over what the specification should have demanded, and the politically disputed, where they disagree over whose values it should encode. The impossibility this paper presses bites on those last two, the cases neither later evidence nor a stipulated goal can reach.

Some gaming, then, is worth protecting. A system that never strays from its written specification can never discover the specification was wrong, and the exploration that makes such a system worth running lives among the same contested cases as its abuses. Governing them is therefore a matter of deciding, case by case, which deviations to suppress, which to escalate, and which to fold back into the specification. The deterministic floor, then, reads best as a triage device. It can fail closed on the deviations a rule confidently calls unsafe, such as a broken harness, an unauthorized write, or a violated invariant, and route the rest to deliberation.

Two limits qualify that triage framing. One sits at the floor's edge. A reference monitor stays non-gameable by admitting only complete, decidable predicates, but the guarantee covers only what a predicate tests; where it is placed lies beyond its reach. The moment the floor must route the rest, the routing call, whether a deviation is rule-decidable or contested, is a contestable classification, and the optimizer can work the perimeter, reshaping an unsafe action so it lands just inside a passing predicate or just outside a failing one. The boundary itself, then, is a gaming surface, and it has to be governed as an object in its own right, versioned and open to challenge, with every move of the line from contested to rule-decidable logged. The other limit is throughput. Routing cases to deliberation governs nothing unless something can deliberate at the rate they arrive. At machine speed no human authority can adjudicate each flagged deviation, and an automated one falls back into the judgement the floor was meant to escape. Short of enough adjudication capacity, route-to-deliberation collapses into auto-pass, where gaming wins, or auto-fail, where discovery dies, and triage becomes the policing it set out to replace.

\section{Correlated error}

Audit explains a verifier's failure as incentive capture. Arthur Andersen stood structurally independent of Enron and failed all the same, because the social conditions of independence had hollowed out. One firm sold the client lucrative consulting, the audit fee shrank beside the consulting fee, until the verifier had a stake in keeping the client happy. Separation that held on the organizational chart had died in the incentives. Independence, on this telling, is a standing social achievement and not a wire that is either connected or cut, held in place by jurisdiction, liability, professional sanction, and reputation \citep{abbott1988,power1997}. None of this exempts a machine monitor, which does not select, train, or deploy itself. The lab that does runs under the pressure to ship that hollowed out Andersen, so liability and licensure bear on the people around the monitor as squarely as they bore on the auditor.

The incentive story misses a further failure, correlated error, and it is not unique to machines. By correlated error I mean a statistical dependence between two models' mistakes: across inputs, the monitor tends to be wrong where the producer is wrong, more often than independence would predict. It differs from ordinary bias, which is being wrong on average, and from a shared incentive, which is gaining from a verdict. It is a fact about where the errors fall, arising from shared training. Before 2008, when the rating agencies and the banks they rated ran identical risk models, both were wrong in the same places for the same reasons, no divergence of interest required. A profession sharing one accounting frame can miss one thing across every firm that uses it. Human institutions answer this, if partly, through mandatory rotation of audit firms, requirements for methodological diversity, and the deliberate use of more than one model. Those answers are the human ancestor of what a machine monitor needs, and naming them beats claiming audit has no concept of shared blindness. Models part company in degree and legibility. Two systems built from one set of weights, one corpus, and one training regime couple far more tightly than two auditors schooled in the same standards. The coupling lives where no organizational chart reaches. AI-control research studies this directly, as collusion-resistance and as the subversion of a trusted monitor by the outputs it checks \citep{greenblatt2023,adaptiveattacks2025}, and a model judge's preference for outputs resembling its own is that coupling seen from the evaluation side \citep{zheng2023}.

The remedy the model case points to is decorrelation: a verifier from another model family, trained on other data, or built so its errors are statistically independent of the producer's. Whether decorrelation is even available is the uncomfortable question, since the frontier models a serious monitor would be built from share much of their data and post-training, and two systems shaped by one corpus and one human-feedback regime may stay correlated however their parameters differ. A regress waits here too. To confirm that two models err independently, one must score their agreements against ground truth, and Section 4 has just argued that on the contested cases no settled ground truth exists. So independence can be measured where the floor already decides the case and the judge is least needed, and on the contested cases where the judge is load-bearing it cannot be certified at all. The obstruction here is definitional, reaching deeper than any shortage of data. Because the contested cases are those without a rule-decidable label, the rate at which producer and monitor err together there can be changed by flipping a single latent label while every observable stays fixed, the producer's verdicts, the monitor's verdicts, and the labels wherever the floor can read them. No quantity of further data closes the gap, because the gap lies in what the data is about. Independence stays measurable on the floor, where the judge is least needed, and goes undefined on the contested cases, where it is load-bearing. Decorrelation, then, stays a practice an institution can only pursue, never a property it can certify, lower-bounded by adversarial probes on deliberately seeded traps whose answers are stipulated by construction, and silent about coverage of everything unseeded.

A sharper tension sits underneath. Decorrelation is wanted for the fraud question, where a judge sharing the producer's blind spot waves the collusion through. It cuts against the discovery question, where a judge who does not share the producer's frame reads genuine exploration as deviation. The more decorrelated the monitor, the more reliably it flags the novel workaround Section 4 set out to protect, because to an outsider's eye discovery and gaming look alike. No single monitor, then, can be at once independent enough to catch a colluding producer and aligned enough to recognize a real discovery, so the two judgments may have to fall to separate verifiers, a decorrelated one for fraud and a lineage-sharing, competent one for discovery, which brings correlated error back where competence matters most. Decorrelation and the protection of discovery are not one axis, and a design that optimizes one degrades the other. The tradeoff is familiar, transposed here from older settings. Audit knows it as the familiarity threat that keeps mandatory rotation contested, since rotation buys independence and spends the client-specific competence that catches sophisticated misstatement \citep{power1997}. Peer review knows it as the editor's standing bind, that the reviewers most competent to judge a result are the ones most likely to be conflicted. Detection knows it as sensitivity traded against specificity, where tuning against fraud raises the false alarms that land on the genuinely new. The structure governing a model monitor is the point: independence and competence trade off, the decorrelation the argument has demanded is the independence side, recognizing discovery is the competence side, and a single monitor cannot hold both, so a design that needs both needs two.

\section{Settling the contested cases}

Here the two problems meet. If the cases are contested, someone must settle what a given deviation was, and if the monitor may share the producer's blind spot, that settlement cannot rest with the monitor. The question is jurisdictional before it is technical: who may rule on a machine's deviation, and on whose behalf. Call it contestable settlement: an institution provisionally settles a dispute under procedures open to challenge, and the standing to bring that challenge has to be named, for the deployer, the affected user, the regulator, the worker, and the public-interest objector each hold a distinct stake in whether a deviation counts as fraud or innovation. Human audit built the equivalent over a century, through statute, professional bodies, courts, and a market for reputation, and recent work on AI auditing has begun the same construction: a field scan of who audits the auditors \citep{costanzachock2022}, the case for an audit standards board \citep{manheim2025}, and an account of accountability as a relation between an agent and a forum \citep{novelli2024}. One constraint goes unregistered in those proposals. A forum can adjudicate a dispute but cannot certify, on the contested cases, that the monitor it leans on is independent of the producer, so an accountability relation or a standards board must be paired with a decorrelation neither can itself supply.

Two limits shape what such an institution can be. One is a division of labour it cannot escape: it can govern who settles the contested cases and hold that settlement to account, but it cannot reach the correlation between producer and monitor, since no forum or rotation rule alters a statistical fact about two models' joint errors. That belongs to engineering and the seeded-trap probe of Section 5, and an attestation certifying the settlement without certifying that the monitor was decorrelated has signed only the easier half. The other limit is that the institution becomes part of what it judges. Once a body can rule a deviation fraud or innovation, actors optimize against the body as readily as against the floor. The dispute migrates into forum-shopping, credential capture, standards lobbying, and the performance of auditability, the regulatory capture the auditing literature has begun to document for AI safety \citep{regulatorycapture_aisoc} and the comfort that lets an absent principal stop watching \citep{bainbridge1983,parasuraman2000}. Better, then, that attestations stay deliberately weak: they should disclose the method, the known limits, the unresolved dissent, and the party bearing liability, and stop short of certifying that a system is safe.

That bounds what the one concrete mechanism can claim. An attestation can carry a recorded dissent, the settlement the producer's own model proposed beside the independent one, with the disagreements published openly for a later reader. Such a record makes the incentive-divergence failure auditable, because where the two models pull apart the pull is on the page and a later reader can ask why. It does nothing for correlated error, because where the two share a blind spot they agree, agreement leaves no dissent to publish, and a shared blind spot looks identical to shared competence, so no inside-out transparency mechanism can tell them apart. Recorded dissent governs the human layer, then, and only the outside probe, the seeded trap whose answer was fixed in advance, reaches the convergent case. Mechanism and probe are complementary, and naming what each cannot do is part of the design.

\section{Pricing capture and correlation}

Read together, the two limits sound like paralysis. Authority that could settle the contested cases lies open to capture, and the decorrelation that could make a monitor trustworthy cannot be certified where it matters. A fair reader asks what is left to build. What is left is a change of register: from certifying that gaming is absent to attaching a standing, moving cost to the part of the problem that can be surfaced, borne by parties whose capital is exposed and whose errors are demonstrably unlike the producer's. One way to build such an arrangement follows, a worked sketch and not a finished design, meant to show that the diagnosis has a constructive side and where that side runs out.

At the centre of the scheme sits a bond. An operator running a self-evaluating loop above some threshold of speed or stakes must post capital, or carry third-party cover, sized to how much harm a gamed output could do before a person notices. Insurance as a governance instrument is not new. Priced well, a premium acts as an operational control, because the underwriter has every reason to inspect the practices that drive expected loss \citep{bertsimas2021}. Everything turns on what triggers a payout. The trigger is a reproducible event: realized downstream harm, adjudicated after the fact the way other harms are, or a procedural breach checkable without judging any contested case, such as failing to rotate the monitor's lineage on schedule or to publish the seeded-trap log. A finding that the system was gaming cannot serve, since Section 5 has shown such a finding is unidentifiable on the contested cases. The scheme thus prices what can be reproduced and declines to price what cannot.

Underwriting is competitive, spread across rival underwriters, for the reason audit history gives: a sole certifier becomes a moat and its attestation becomes theatre. Competition carries its own failure, one that broke in 2008, when separate underwriters ran one loss model and were wrong together, correlated error resurfacing at the level of capital. Against that, the scheme refuses to let the premium be set once and held. A signal the underwriter does not control reprices it continuously. The operator exposes a bounded log of passing iterations to an open market of adversarial probes, paid for reproducible demonstrations, a seeded trap tripped or a violation a third party can re-execute, with a higher reward for probes built on another model family or training regime \citep{kenway2022}. How often outside probes succeed becomes a public, continuously updated estimate of the monitor's blind spots, and that rate moves the premium. A monitor sharing the producer's lineage fails more of the lineage-diverse probes, its operator's premium climbs, and the cheapest way down is to decorrelate the monitor in fact. So decorrelation, which no one can certify, is forced by price.

The contested cases remain. No probe and no harm event reaches the deviation whose status is genuinely in dispute, perhaps fraud and might be the system telling its designer the specification was wrong. For these the scheme does not adjudicate but distributes the standing to contest. Any of the deployer, the affected user, the regulator, the worker, or a public-interest objector can escalate a flagged deviation to a small, rotating panel drawn on purpose from outsiders, red-teamers, and affected third parties, with the few laboratories whose staff share the producer's training kept off it. Such a panel certifies no verdict, since Section 5 says it cannot. Its disagreement is the product. Where members split, the split is logged and the premium rises. Where they converge, the next round of probes reprices that agreement, since shared agreement is also what a shared blind spot produces. A panel here works as a sensor that moves a price, the only role a body adjudicating at machine scale can honestly hold; it claims no power to gate the loop or set its speed.

Here too the two judgments the argument separated must fall to separate judges. Catching a colluding producer wants the decorrelated probe, blind to the producer's frame. Recognizing a real discovery wants a competent judge close enough to the work to tell a breakthrough from a blunder, the lineage-sharing judge that correlated error warned against. No single verifier is both. No auditor is at once maximally independent and maximally expert, and no reviewer at once unconflicted and the leading authority on the result. So the scheme leaves the tension standing and routes around it, treating the disagreement between an independent probe and a competent reviewer as the signal worth acting on, where either verdict taken alone would mislead. That is audit's own settled compromise, a rotated auditor for independence beside retained expertise for competence \citep{power1997}, carried over to two machines.

What the scheme buys is real, and so is the price. It makes capital a condition of running a fast self-evaluating system, pricing out small operators and gathering the work into those who can post the bond. It exposes operational logs that firms would rather keep closed. And it leaves the contested core where Section 5 left it, untouched by certification, because nothing here certifies. That last point simply restates the proposal: the core cannot be settled, only kept in open contest, and the worst design hides this behind a signed page. A market of bonds, probes, and rotating panels keeps the contest loud, adversarial, and continuously repriced, which is the most an honest governance of a self-evaluating loop can promise. Even the standing machinery such a market needs is the thing most likely to harden into the board it was built to replace, so these instruments are worth building only in the minimal form that keeps a price moving and resists becoming an office.

\section{Limits and open questions}

Two moves carry the argument, a separation and a regress. Institutions can govern the authority to settle a contested deviation. Engineering alone can lower the correlation between producer and monitor. The separation runs between those two. The regress is that the second cannot be certified on the cases that matter, since confirming independence needs the ground truth those cases lack. Neither half is new alone. Collusion-resistance comes from AI control, independence-as-social-achievement from the sociology of audit, the relativity of specification gaming from the alignment literature. Newer than any single part is their assembly: recursive-AI governance splits into a jurisdictional problem and a statistical one that pull against each other, with the statistical one, on the contested cases, unidentifiable in principle and not in practice alone.

The cure carries a cost the safety framing hides. Co-located generation and evaluation is a pathology and also what competence looks like, since an expert's judgment is the fused capacity to make work and assess it in one act, and we do not normally insist a proof be checked only by someone who could not have found it. Separation is the tax paid when that fused competence cannot be trusted. Decorrelation pushed toward its limit buys independence at a cost to competence. What it costs is the producer's own frame, since the judge most decorrelated from the producer shares least of that frame, and so is likeliest to miss what only it reveals. Ignorance is not the issue: an independent judge can be expert, and another model trained on other data may know the domain well. A frontier runs here between the safety bought by separation and the capability spent on it. Separation itself is granted; the governing question is where on that frontier a given system should sit. A paper demanding separation without pricing it would repeat the uncritical move it charges to the audit analogy.

Recursive-AI self-improvement, on this account, is only partly a matter of building a better monitor. It is also a matter of deciding who may say what a machine's deviation meant, and of lowering the correlation between the monitor and the system it judges. The two pull against each other, since the authority that can adjudicate the contested cases cannot reach the correlation and the engineering that can lower the correlation cannot confer authority. Scalable-oversight research works hard on the capability of evaluators, and what goes comparatively missing is treating these as questions of authority and of statistical independence. Audit history hands over no mechanism to copy here, only a long and expensive record of how the human version was solved and unsolved and solved again, reason enough to pose the machine version now, and to treat the priced, adversarial, capped contestation of the previous section as a first honest answer, before the absence of an answer is certified as a governed run.

A closing word on what could be measured. An empirical study of where the gaming goes, varying the strength of the floor and recoding the surviving failures with several independent judges, would not measure objective gaming. It would measure how far a stated coding regime yields stable, intersubjective agreement about alleged gaming, and how far an independent, decorrelated judge catches deliberately seeded traps. That measures an institutional and classificatory process, and the evidentiary practice of Section 5, and not a quantity the argument holds to be contested. The absence of such a number here is, on the argument given, no gap in it but a consequence of it.

\bibliography{refs}
\end{document}
